\chapter{Settings and Preliminaries}\label{ch:examples}

This chapter establishes the foundational framework for the free boundary min-max theory in complete Riemannian manifolds.

\section{Notation and Basic Conventions}\label{sec:notation}

\subsection{Ambient Manifold and Doubled Manifold}\label{subsec:ambient-manifold}

Let $(M^{n+1}, g)$ be a smooth, complete Riemannian manifold of dimension $n+1$ with non-empty boundary $\partial M$. We assume throughout that $\partial M$ is a smooth, closed, $n$-dimensional submanifold of $M$. The metric $g$ induces a Riemannian metric on $\partial M$ by restriction, which we also denote by $g$ when no confusion arises.

By Nash's isometric embedding theorem~\cite{Nash56}, we may assume that $M$ is isometrically embedded in some Euclidean space $\mathbb{R}^N$ with $N$ sufficiently large. When working with varifolds and currents, this ambient Euclidean structure is often convenient, though all our results are intrinsic to $(M, g)$.

A fundamental construction in free boundary min-max theory is the \emph{doubled manifold}. Let
$\widetilde{M}$ denote the manifold obtained by reflecting $M$ across $\partial M$ and gluing the two
copies smoothly. More precisely, $\widetilde{M}$ is the disjoint union $M \sqcup M^-$ (where $M^-$ is
a copy of $M$) with the boundaries $\partial M$ and $\partial M^-$ identified, and equipped with a
smooth Riemannian metric that agrees with $g$ on $M$ and its reflection on $M^-$. The doubled manifold
$\widetilde{M}$ is a complete Riemannian manifold without boundary, and we identify $M$ with one of the
two copies in $\widetilde{M}$.

\begin{notation}
For any subset $U \subset M$, we denote by $\widetilde{U}$ the corresponding open set in $\widetilde{M}$
obtained by reflecting $U$ across $\partial M$. If $U \cap \partial M = \emptyset$, then
$\widetilde{U} = U \sqcup U^-$ is the disjoint union of two copies. If $U \cap \partial M \neq \emptyset$,
the reflection is performed in Fermi coordinates near $\partial M$ to ensure smoothness.
\xinze{Smoothness should be guaranteed no other clarification needed.}
\end{notation}

We denote by $\nu$ the outward-pointing unit normal vector field on $\partial M$ (pointing out of $M$
into $M^-$ in $\widetilde{M}$). For $p \in \partial M$, the tangent space $T_p M$ decomposes as
$T_p M = T_p(\partial M) \oplus \mathbb{R} \nu(p)$.

\subsection{Measures, Norms, and Currents}\label{subsec:measures-norms}
\xinze{Is this section used anywhere?}
We denote by $\mathcal{H}^k$ the $k$-dimensional Hausdorff measure on $M$ (or $\widetilde{M}$). For a smooth $k$-dimensional submanifold $\Sigma \subset M$,
the measure $\mathcal{H}^k \llcorner \Sigma$ coincides with the $k$-dimensional volume measure on
$\Sigma$ induced by the Riemannian metric. 

For integral currents, we denote by $\mathbf{M}(T)$ the \emph{mass} of a current $T$, and by
$\mathcal{F}(T)$ the \emph{flat norm}. Recall that for a $k$-dimensional integral current
$T \in \mathcal{I}_k(M)$, the mass is defined by
\begin{equation*}
    \mathbf{M}(T) := \sup\left\{ T(\omega) : \omega \in \Omega^k(M), \|\omega\|_{C^0} \leq 1 \right\},
\end{equation*}
where $\Omega^k(M)$ denotes the space of smooth $k$-forms on $M$. A $k$-dimensional integral current in $M$ with \emph{free boundary} is an element of $\mathcal{Z}_k(M, \partial M; \mathbb{Z}_2)$, the space of integral $k$-cycles modulo $2$ with boundary supported on $\partial M$. More precisely, $T \in \mathcal{Z}_k(M, \partial M; \mathbb{Z}_2)$ if $T$ is an integral current with $\partial T \in \mathcal{I}_{k-1}(\partial M; \mathbb{Z}_2)$.

When a hypersurface $\Gamma \subset M$ is the reduced boundary of a Caccioppoli set $\Omega \subset M$ (see Section~\ref{sec:caccioppoli}), we write $\Gamma = \bdy \Omega$ in our unified notation. The current $[\Gamma]$ associated to $\Gamma$ (with an
appropriate choice of orientation or $\mathbb{Z}_2$ coefficients) satisfies
$\mathbf{M}([\Gamma]) = \mathcal{H}^n(\Gamma) = \Per_M(\Omega)$.

\subsection{Vector Fields and Tangent Spaces}\label{subsec:vector-fields}

Let $\mathfrak{X}(M)$ denote the space of smooth vector fields on $M$ with compact support. When
working with the doubled manifold, we also consider $\mathfrak{X}(\widetilde{M})$, the space of smooth
vector fields on $\widetilde{M}$ with compact support. For free boundary variations, we often restrict
to vector fields that are \emph{tangent} to $\partial M$: we denote by $\mathfrak{X}_{\tan}(M)$ the
space of smooth vector fields $X$ on $M$ with compact support such that $X(p) \in T_p(\partial M)$ for
all $p \in \partial M$. Similarly, $\mathfrak{X}_{\perp}(M)$ denotes vector fields that are
\emph{normal} to $\partial M$, i.e., $X(p) \in \mathbb{R} \nu(p)$ for all $p \in \partial M$.

Any vector field $X \in \mathfrak{X}(M)$ decomposes near $\partial M$ as $X = X_{\tan} + X_{\perp}$,
where $X_{\tan} \in \mathfrak{X}_{\tan}(M)$ and $X_{\perp} \in \mathfrak{X}_{\perp}(M)$. This
decomposition plays a fundamental role in the first variation formula for free boundary varifolds
(Section~\ref{sec:stationary-varifolds}).

For a vector field $X \in \mathfrak{X}(M)$, we denote by $\dive X$ its divergence with respect to the
Riemannian metric $g$. In local coordinates $(x^1, \ldots, x^{n+1})$, if $X = X^i \partial_i$ and
$g = g_{ij} dx^i dx^j$, then
\[
    \dive X = \frac{1}{\sqrt{\det g}} \partial_i\left( \sqrt{\det g} \, X^i \right).
\]
The divergence theorem in the relative setting will be used extensively in the coarea formula
(Section~\ref{sec:coarea}) and in variation formulas.

\subsection{Sweepouts: Preliminary Definition}\label{subsec:sweepouts-prelim}

A \emph{sweepout} of a relative open set $U \subset M$ is, roughly speaking, a continuous family
$\{\Gamma_t\}_{t \in [0,1]}$ of hypersurfaces that ``fills'' $U$ as $t$ varies from $0$ to $1$, with
$\Gamma_0$ and $\Gamma_1$ being ``trivial'' (e.g., empty or outside of $U$). The precise definition, including regularity, continuity, and mass bounds, is given in Definition~\ref{def:sweepout}, as it requires the notion of families of surfaces developed there.
For a technical variant used in the appendix for interpolation lemmas, see also
Definition~\ref{def:sweepout-appendix} in Appendix~\ref{ch:interpolation-lemma}.

For now, we introduce the \emph{height} functional: for a family $\{\Gamma_t\}_{t \in [0,1]}$, we define
\[
    \mathcal{F}(\{\Gamma_t\}) := \sup_{t \in [0,1]} \mathcal{H}^n(\Gamma_t \cap \mathrm{int}(M)),
\]
where $\mathrm{int}(M) = M \setminus \partial M$ denotes the interior of $M$. The height
$\mathcal{F}(\{\Gamma_t\})$ measures the ``cost'' of the sweepout in terms of maximal area. In min-max
theory, we seek to minimize this height over all sweepouts in a given homotopy class, yielding the
\emph{width} (see Section~\ref{sec:good-sets}).

\section{Caccioppoli Sets}\label{sec:caccioppoli}

In this section, we introduce the theory of Caccioppoli sets.

\subsection{Perimeter and Reduced Boundary}\label{subsec:perimeter}
\begin{definition}[Caccioppoli set]\label{def:caccioppoli} Let $M$ be a complete Riemannian manifold (possibly with boundary). A measurable set $E \subset M$ is
called a \emph{Caccioppoli set} (or a set of \emph{finite perimeter}) if its \emph{perimeter} in $M$,
defined by
\begin{equation*}
    \Per_M(E) := \sup\left\{ \int_E \dive \omega\, dV : \omega \in \mathfrak{X}(M),\, \|\omega\|_{C^0(M)} \leq 1 \right\},
\end{equation*}
is finite.
\end{definition}

\begin{definition}[Reduced boundary]\label{def:reduced-boundary}
Let $E \subset M$ be a Caccioppoli set. The \emph{reduced boundary} $\partial^* E$ is the set of points $x \in \spt{\mu_E}$ such that the limit 
\begin{equation*}
    \nu_E(x) := \lim_{r \to 0^+} \frac{\int_{B_r(x)} \nabla \chi_E\, dV}{|\int_{B_r(x)} \nabla \chi_E\, dV|}
\end{equation*}
exists and belongs to $S^{n - 1}$. We call the vector field $\nu_E : \partial^* E \to TM$ the \emph{(measure-theoretic) outward unit normal} to $E$.
\end{definition}

The reduced boundary has a rich geometric structure, as described by the following fundamental theorem:

\begin{theorem}[Structure of reduced boundary, {\cite[Theorem~15.9]{Mag12}}]\label{thm:structure-reduced-boundary}
Let $E \subset M$ be a Caccioppoli set. Then:
\begin{enumerate}[label=(\roman*)]
    \item The reduced boundary $\partial^* E$ is countably $n$-rectifiable, i.e., there exist countably
    many Lipschitz maps $f_i : \mathbb{R}^n \to M$ such that
    \[
        \mathcal{H}^n\left( \partial^* E \setminus \bigcup_{i=1}^\infty f_i(\mathbb{R}^n) \right) = 0.
    \]
    \item The perimeter of $E$ is given by the Hausdorff measure of the reduced boundary:
    \[
        \Per_M(E) = \mathcal{H}^n(\partial^* E \cap M).
    \]
    \item The map $x \mapsto \nu_E(x)$ is $\mathcal{H}^n$-measurable on $\partial^* E$, and for any
    Borel set $A \subset M$ and any $\omega \in \mathfrak{X}(M)$,
    \[
        \int_{E \cap A} \dive \omega\, dV = \int_{\partial^* E \cap A} \langle \omega, \nu_E \rangle\, d\mathcal{H}^n.
    \]
\end{enumerate}
\end{theorem}

Part (iii) of the theorem is a generalization of the divergence theorem to Caccioppoli sets. It shows
that the perimeter functional can be computed as an integral over the reduced boundary, just as in the
smooth case.

\subsection{Integral Currents}\label{subsec:caccioppoli-currents}

\begin{theorem}[Cycles, {\cite[Section~4.5]{Mag12}}]\label{thm:caccioppoli-cycles}
Let $E \subset M$ be a Caccioppoli set with $\Per_M(E) < \infty$. Then the boundary of $E$ (in the
sense of currents) can be identified with an integral $n$-current modulo $2$:
\[
    \partial E = \nu_E \, \mathcal{H}^n \llcorner \partial^* E \in \mathcal{Z}_n(M; \mathbb{Z}_2),
\]
where $\mathcal{Z}_n(M; \mathbb{Z}_2)$ denotes the space of integral $n$-cycles with $\mathbb{Z}_2$
coefficients. The mass of this cycle coincides with the perimeter:
\[
    \mathbf{M}(\partial E) = \Per_M(E) = \mathcal{H}^n(\partial^* E).
\]
\end{theorem}

\subsection{Perimeter in the Free Boundary Setting}\label{subsec:perimeter-free-boundary}

When working with manifolds with boundary, we must adapt the notion of perimeter to account for the presence of $\partial M$. For a Caccioppoli set $E \subset M$, the reduced boundary $\partial^* E$ may intersect $\partial M$. The portion of $\partial^* E$ on $\partial M$ corresponds to the \emph{free boundary} of $E$, while the portion in the interior of $M$ corresponds to the \emph{interior boundary}.

\begin{definition}[Relative perimeter]\label{def:relative-perimeter}
Let $U \subset M$ be a relatively open set, and let $E \subset M$ be a Caccioppoli set. The \emph{relative perimeter} of $E$ in $U$ is defined by
\[
    \Per_M(E, U) := \sup\left\{ \int_{E \cap U} \dive \omega\, dV : \omega \in \mathfrak{X}(M),\,
    \mathrm{supp}(\omega) \subset U,\, \|\omega\|_{C^0} \leq 1 \right\}.
\]
This measures the ``boundary area'' of $E$ localized to the region $U$.
\end{definition}
The relative perimeter satisfies a localized version of Theorem~\ref{thm:structure-reduced-boundary}:
\[
    \Per_M(E, U) = \mathcal{H}^n(\partial^* E \cap U).
\]
This quantity plays a crucial role in the non-excessive min-max theory, where we control the perimeter
of slices localized to small balls.

\begin{notation}
When the ambient manifold $M$ is clear from context, we write $\Per(E)$ instead of $\Per_M(E)$, and
$\Per(E, U)$ instead of $\Per_M(E, U)$. For a family of Caccioppoli sets $\{E_t\}_{t \in [0,1]}$, we
denote by $\{\Gamma_t\}$ the corresponding family of reduced boundaries $\Gamma_t = \partial^* E_t$,
and we identify $\Gamma_t$ with the current $\partial E_t \in \mathcal{Z}_n(M; \mathbb{Z}_2)$ via
Theorem~\ref{thm:caccioppoli-cycles}.
\end{notation}

\section{Varifolds and Rectifiability}\label{sec:varifolds}
\subsection{Rectifiable and Integral Varifolds}\label{subsec:rectifiable-varifolds}
We recall the basic theory for rectifiable and integral varifold; see \cite{Sim84} for more information. If $U$ is an open subset of $M$, any Radon measure on Grassmannian $G(U)$ of unoriented $n$-planes on $U$ is said to be an $n$-\textit{varifold} in $U$. We denote the space of $n$-varifolds by $\mathcal{V}(U)$ and the topology it endowed with is weak*-topology. Thus, a sequence ${V^k} \subset \mathcal{V}(U)$ converges to $V$ is 
\begin{equation*}
    \lim_{k \to \infty}\int \varphi(x, \pi)dV^k(x, \pi) = \int \varphi(x, \pi)dV(x, \pi)
\end{equation*}
for every $\varphi \in C_c(G(U))$. Here $\pi$ denotes an $n$-plane of $T_xM$. If $U' \subset U$ and $V \in \mathcal{V}(U)$, then $V \mres U'$ is the restriction of the measure $V$ to $G(U')$.

Moreover, $\norm{V}$ is the nonnegative measure on $U$ defined by
\begin{equation*}
    \int_U \varphi(x)d\norm{V}(x) = \int_{G(U)}\varphi dV(x, \pi), \quad \forall \varphi \in C_c(U).
\end{equation*}
The support of $\norm{V}$, denoted by $\spt{\norm{V}}$, is the smallest closed set outside which $\norm{V}$ vanishes identically. The number $\norm{V}(U)$ will be called the \textit{mass} if $V$ in $U$.

Recall also that an $n$-dimensional rectifiable set is the countable union of closed subsets of $C^1$ surfaces (modulo sets of $\haus^n$-measure $0$). If $R \subset U$ is an $n$-dimensional rectifiable set and $h: R \to \R^+$ is a Borel function, then the \textit{varifold $V$ induced by $R$} is defined by
\begin{equation*}
    \int_{G(U)}\varphi(x, \pi)dV(x, \pi) = \int_R h(x) \varphi(x, T_x R)d\haus^n(x)
\end{equation*}
for all $\varphi \in C_c(G(U))$. Here $T_xR$ denotes the tangent plane to $R$ in $x$. If $h$ is integer-valued, then we say that $V$ is an \textit{integral rectifiable varifold}. If $\Sigma = \bigcup n_i \Sigma_i$, then we say that $V$ is an \textit{integral rectifiable varifold.}

\section{Stationary Varifolds with Free Boundary}\label{sec:stationary-varifolds}
In this section, we develop the variational theory of varifolds in manifolds with boundary, focusing on the free boundary setting relevant to our min-max constructions. The key concept is that of a \emph{stationary varifold}, which generalizes the notion of minimal hypersurface to the measure-theoretic framework.

\subsection{First Variation and Stationarity}\label{subsec:first-variation}

Let $(M^{n+1}, g)$ be a complete Riemannian manifold with boundary $\partial M$, and let $\widetilde{M}$
denote the doubled manifold (Section~\ref{subsec:ambient-manifold}). For a varifold
$V \in \mathcal{V}_n(M)$, the \emph{first variation} of $V$ with respect to a vector field
$X \in \mathfrak{X}(\widetilde{M})$ measures the rate of change of mass under the flow generated by $X$.

\begin{definition}[First variation of a varifold]\label{def:first-variation-varifold}
Let $V \in \mathcal{V}_n(\widetilde{M})$ be a varifold. For a smooth vector field
$X \in \mathfrak{X}(\widetilde{M})$ with compact support, the \emph{first variation} of $V$ with respect
to $X$ is defined by
\[
    \delta V(X) := \int_{G_n(\widetilde{M})} \operatorname{div}_S X(x)\, dV(x, S),
\]
where $\operatorname{div}_S X(x)$ denotes the tangential divergence of $X$ at $x$ along the $n$-plane
$S \in G_n(T_x \widetilde{M})$. Explicitly, if $\{e_1, \ldots, e_n\}$ is an orthonormal basis for $S$,
then
\[
    \operatorname{div}_S X(x) := \sum_{i=1}^n \langle \nabla_{e_i} X, e_i \rangle,
\]
where $\nabla$ is the Levi-Civita connection on $\widetilde{M}$.
\end{definition}

For a smooth hypersurface $\Sigma \subset M$ with associated varifold $V = [\Sigma]$, the first
variation has a geometric interpretation: if $X$ is a vector field and $\phi_t$ is the flow generated
by $X$, then
\[
    \delta V(X) = \left. \frac{d}{dt} \right|_{t=0} \mathcal{H}^n(\phi_t(\Sigma)).
\]
Using the divergence theorem, this can be expressed in terms of the mean curvature vector $H$ of $\Sigma$:
\[
    \delta V(X) = -\int_\Sigma \langle H, X \rangle\, d\mathcal{H}^n + \int_{\partial \Sigma} \langle \nu_\Sigma, X \rangle\, d\mathcal{H}^{n-1},
\]

\begin{definition}[Stationary varifold with free boundary]\label{def:stationary-free-boundary}
Let $V \in \mathcal{V}_n(M)$ be a varifold. We say that $V$ is \emph{stationary with free boundary} if:
\begin{enumerate}[label=(\roman*)]
    \item $\delta V(X) = 0$ for all $X \in \mathfrak{X}_{\tan}(M)$, i.e., for all vector fields tangent
    to $\partial M$;
    \item $\|V\|(\partial M) = 0$, i.e., the varifold has no mass on the boundary itself.
\end{enumerate}
\end{definition}

\begin{remark}
Condition (i) ensures that $V$ is a critical point for area variations that preserve the boundary
$\partial M$. For a smooth hypersurface $\Sigma$ with $\partial \Sigma \subset \partial M$, this is
equivalent to requiring:
\begin{itemize}
    \item $H = 0$ on $\Sigma \cap \mathrm{int}(M)$ (zero mean curvature in the interior);
    \item $\Sigma$ meets $\partial M$ orthogonally along $\partial \Sigma$ (free boundary condition).
\end{itemize}
Condition (ii) is a technical requirement ensuring that the varifold does not concentrate mass on the
boundary.
\end{remark}

\subsection{Monotonicity Formula}\label{subsec:monotonicity}

One of the fundamental tools in the analysis of stationary varifolds is the monotonicity formula, which provides control on the density of the varifold at different scales. For free boundary varifolds, the monotonicity formula must be adapted to account for the presence of $\partial M$.
\begin{theorem}[Monotonicity formula for free boundary varifolds~\cite{LZ17}]\label{thm:monotonicity-free-boundary}
Let $V \in \mathcal{V}_n(M)$ be a stationary varifold with free boundary (Definition~\ref{def:stationary-free-boundary}). Assume that $\partial M$ is smooth. Then for any $x_0 \in M$ and any $0 < r < R < \mathrm{dist}(x_0, \infty)$ (where $\mathrm{dist}(x_0, \infty) = \infty$ if $M$ is compact), we have
\[
    \frac{\|V\|(B_R(x_0))}{R^n} \geq \frac{\|V\|(B_r(x_0))}{r^n}.
\]
Moreover, if $x_0 \in \mathrm{int}(M)$ and $R < \mathrm{dist}(x_0, \partial M)$, then the ratio
$\|V\|(B_r(x_0))/r^n$ is non-decreasing in $r$ for $r \in (0, R)$.
\end{theorem}

\xinze{when the following density bound is used?}
\begin{corollary}[Density bounds]\label{cor:density-bounds}
Let $V \in \mathcal{V}_n(M)$ be a stationary varifold with free boundary. Then the density
$\Theta^n(V, x)$ exists for $\|V\|$-almost every $x \in M$ and
satisfies:
\begin{enumerate}[label=(\roman*)]
    \item $\Theta^n(V, x) \geq \frac{1}{2}$ for all $x \in \mathrm{spt}(V) \cap \mathrm{int}(M)$;
    \item $\Theta^n(V, x) \geq \frac{1}{4}$ for all $x \in \mathrm{spt}(V) \cap \partial M$ (assuming
    $\partial M$ is smooth).
\end{enumerate}
\end{corollary}

\subsection{Mean Convexity and Maximum Principle}\label{subsec:max-principle}

In the free boundary setting, the geometry of $\partial M$ plays an important role in determining the
behavior of stationary varifolds. A key condition is \emph{mean convexity}.

\xinze{where mean convexity is used?}
\begin{definition}[Mean convexity]\label{def:mean-convexity}
Let $(M, g)$ be a Riemannian manifold with boundary $\partial M$. We say that $\partial M$ is \emph{mean convex} if the mean curvature $H_{\partial M}$ of $\partial M$ (computed with respect to the inward-pointing normal) is non-negative at every point. If $H_{\partial M} > 0$ everywhere, we say $\partial M$ is \emph{strictly mean convex}.
\end{definition}

\begin{remark}
For example, if $M = B_R(0) \subset \mathbb{R}^{n+1}$ is a Euclidean ball, then $\partial M = \partial B_R(0)$ is strictly mean convex with $H_{\partial M} = n/R > 0$.
\end{remark}

Mean convexity of $\partial M$ has important consequences for stationary varifolds:
\xinze{We may not need mean convexity in maximum principle}
\begin{theorem}[Maximum principle for free boundary varifolds]\label{thm:max-principle}
Let $V \in \mathcal{V}_n(M)$ be a stationary varifold with free boundary, and assume $\partial M$ is
strictly mean convex. Then:
\begin{enumerate}[label=(\roman*)]
    \item If $\Sigma \subset \mathrm{spt}(V)$ is a smooth component with $\partial \Sigma \subset \partial M$, then $\Sigma$ cannot have an interior tangency with $\partial M$. That is, if $x \in \Sigma \cap \partial M$ and $x \notin \partial \Sigma$, then we must have $x \in \partial \Sigma$ (contradiction), i.e., no false boundary points exist.
    \item More generally, $\mathrm{spt}(V) \cap \partial M$ has codimension at least $1$ in $\partial M$.
\end{enumerate}
\end{theorem}
\subsection{Second Variation, Stability, and Morse Index for Varifolds}\label{subsec:varifold-second-variation}

\xinze{New subsection: Added to define varifold second variation and Morse index before regularity, resolving circular reasoning in Chapter 8.}

For stationary varifolds, the second variation of area and the associated Morse index can be defined even when the varifold has singularities. This is crucial for min-max theory, where one needs to establish stability or index bounds \emph{before} proving regularity. 

\begin{definition}[Second variation for varifolds]\label{def:second-variation-varifold}
Let $V \in \mathcal{V}_n(M)$ be a stationary integral varifold with free boundary. For a compactly
supported normal variation $\phi \in C^1_c(\mathrm{reg}\, V; \nu)$ (i.e., $\phi$ is supported in the
regular part $\mathrm{reg}\, V$ and takes values in the normal bundle $\nu$ to $V$), the
\emph{second variation} is defined by:
\[
    \delta^2 V(\phi, \phi) = \int_{\mathrm{reg}\, V} \left( |\nabla^\perp \phi|^2 - |A|^2 |\phi|^2 - \mathrm{Ric}(\nu, \nu)|\phi|^2 \right) d\|V\|,
\]
where:
\begin{itemize}
    \item $\nabla^\perp$ denotes the covariant derivative in the normal bundle;
    \item $A$ is the second fundamental form of $V$ (defined $\|V\|$-a.e.\ on $\mathrm{reg}\, V$);
    \item $\mathrm{Ric}$ is the Ricci curvature of the ambient manifold $(M, g)$;
    \item $\nu$ is a choice of unit normal to $V$ (which exists locally on $\mathrm{reg}\, V$ when
    $V$ is orientable).
\end{itemize}
\end{definition}

\begin{definition}[Stability and Morse index for varifolds]\label{def:morse-index-varifold}
Let $V \in \mathcal{V}_n(M)$ be a stationary integral varifold with free boundary. We say that:
\begin{enumerate}[label=(\roman*)]
    \item $V$ is \emph{stable} if for each open set $\Omega \subset M$ with $\mathrm{sing}\, V \cap \Omega = \emptyset$
    (when $2 \leq n \leq 6$), the second variation satisfies $\delta^2 V(\phi, \phi) \geq 0$ for all
    normal variations $\phi \in C^1_c(\mathrm{reg}\, V \cap \Omega; \nu)$.

    \item The \emph{Morse index} (or simply \emph{index}) of $V$ is defined by:
    \[
        \mathrm{Index}(V) := \max\left\{ \dim W : W \subset C^1_c(\mathrm{reg}\, V; \nu), \, \delta^2 V(\phi, \phi) < 0 \, \forall \phi \in W \setminus \{0\} \right\}.
    \]
    In other words, $\mathrm{Index}(V)$ is the maximal dimension of a subspace on which the second
    variation is negative definite.
\end{enumerate}
\end{definition}
\xinze{Give the reference or give the proof.}
\subsection{Reflection Principle}\label{subsec:reflection}
\begin{lemma}[Reflection principle for free boundary varifolds]\label{lem:reflection-principle}
Let $V \in \mathcal{V}_n(M)$ be a stationary varifold with free boundary. Assume $\partial M$ is smooth (at least $C^2$). Then $V$ can be extended to a stationary varifold $\widetilde{V} \in \mathcal{V}_n(\widetilde{M})$ in the doubled manifold $\widetilde{M}$ (Section~\ref{subsec:ambient-manifold}) by reflection across $\partial M$. Moreover, if $V$ is integral and has multiplicity one, then so does $\widetilde{V}$.
\end{lemma}

\section{Coarea Formula and Isoperimetric Inequalities}\label{sec:coarea}
The coarea formula and isoperimetric inequalities are fundamental tools in geometric measure theory that allow us to control the size of level sets and boundaries in terms of ambient volumes. We present both the standard versions (without boundary) and the relative versions adapted to manifolds with boundary. These tools will be used repeatedly throughout the thesis to estimate areas of slices and to control the mass of sweepouts.

\subsection{Coarea Formula}
We first recall the classical coarea formula, which relates the integral of a function to the measures of its level sets.
\xinze{provide refernces for classical version here}
For manifolds with boundary, we work in the doubled manifold $\widetilde M$ to avoid boundary
complications. This leads to the following relative version:

\begin{theorem}[Relative coarea formula]\label{thm:coarea-formula}
Let $N\subset \widetilde M$ be a smooth $(n+1)$-dimensional submanifold whose boundary is tangent to
$\partial M$ along $\partial N\cap \partial M$. Let $f:N\to \mathbb{R}$ be a Lipschitz function and
denote by $N_t:=f^{-1}(t)$. Then, for any Borel set $A\subset N$,
\[
    \int_A |\nabla f|\,d\mathcal{H}^{n+1}
    =\int_{\mathbb{R}} \mathcal{H}^n(A\cap N_t)\, dt.
\]
In particular, if $f$ is proper on $N$ and has no critical points in $A$, then for almost every $t$
the slice $N_t\cap M$ is a smooth hypersurface whose area is controlled by the integral on the left.
\end{theorem}

\subsection{Isoperimetric Inequalities}
\begin{theorem}[Relative isoperimetric inequality]\label{thm:relative-isoperimetric}
Let $M$ be a manifold with boundary $\partial M$, and let $U\subset M$ be a relative open set with
finite volume. Let $x\in M$ and suppose $B_r(x)$ is a geodesic ball of radius $r>0$ that is
$(1+L)$-bi-Lipschitz diffeomorphic to a Euclidean ball (or half-ball if $x\in\partial M$). Then
for any finite-perimeter set $\Omega\subset U$, there exists a radius $s\in(0,r)$ such that
\begin{equation*}
    \mathcal{H}^n(\partial\Omega\cap \partial B_s(x)\cap M)
    \leq \frac{C(n)(1+L)^n}{r}\,\mathcal{H}^{n+1}(\Omega\cap B_r(x)\cap M)
\end{equation*}
where $C(n)>0$ is a constant depending only on the dimension.
\end{theorem}

\begin{remark}
In applications, we often take $L=L_{\mathrm{ball}}$ where $(1+L_{\mathrm{ball}})^n\le 1+\varepsilon$
for some small $\varepsilon>0$, so that the bi-Lipschitz constant can be absorbed into the error
term. This choice appears frequently in Chapters~\ref{ch:nested-sweepout} and~\ref{ch:no-escape-to-infinity}.
\end{remark}

\section{Regularity Theory for Stationary Varifolds}\label{sec:regularity-theory}
Once a stationary integral varifold is known to be integral and rectifiable, one can apply powerful regularity theorems to conclude smoothness. The fundamental result in this direction for codimension 1 varifolds is due to Wickramasekera~\cite{Wic14}, which we now summarize. This section collects the main regularity results that will be used in the min-max construction developed in Chapter~\ref{ch:alternative-regularity}.

\begin{theorem}[Wickramasekera~\cite{Wic14}]\label{thm:wickramasekera-prelim}
Let $V$ be a stable integral $n$-varifold in a smooth $(n+1)$-dimensional Riemannian manifold with
$2 \leq n \leq 6$. Suppose $V$ satisfies the following \emph{$\alpha$-Structural Hypothesis} for
some $\alpha \in (0,1/2)$: no singular point has a neighborhood in which the support of $V$
corresponds to a union of three or more embedded $C^{1,\alpha}$ hypersurfaces-with-boundary meeting
only along their common boundary.

Then $\mathrm{sing}\, V = \emptyset$, i.e., $V$ corresponds to a smooth embedded hypersurface.
\end{theorem}
\xinze{Explain what is $\alpha$-Structural Hypothesis}
The $\alpha$-Structural Hypothesis is automatically satisfied whenever the singular set has vanishing
$(n-1)$-dimensional Hausdorff measure. In particular, any integral varifold that arises as a limit
of currents without cancellation (such as our multiplicity-one case in Chapter~\ref{ch:alternative-regularity})
satisfies this hypothesis. The theorem also applies in higher dimensions with optimal estimates on
the singular set; see Appendix~\ref{sec:wickramasekera-detailed} for a more complete statement.

\chapter{Families of Surfaces and Sweepouts}\label{ch:sweepouts}

We begin by defining what it means for a family of hypersurfaces to vary continuously especially in the setting of manifold with boundary. 
\begin{definition}[Family of hypersurfaces]\label{def:family-surfaces}
A family $\br{\Gamma_t}$, $t \in [0, 1]$, of closed subsets $M$ with finite Hausdorff measure will be called a family of hypersurfaces if:
\begin{itemize}
    \item[(s1)] for each $t$, there is a finite $P_t \subset \text{int}(M)$ such that $\Gamma_t\backslash \partial M$ is a smooth hypersurfaces in $M \backslash P_t$;
    \item[(s2)] $\haus^n(\Gamma_t)$ depends smoothly on $t$ and $t \to \Gamma_t$ is continuous in the Hausdorff sense \xinze{This is true in general as we can perturb surfaces. or we can assume the boundary being smooth}
    \item[(s3)] on any $U \cptsub \text{int}(M)\backslash P_{t_0}$, $\Gamma_t\to\Gamma_{t_0}$ smoothly in $U$ as $t \to t_0$;
    \item[(s4)] (no concentration of mass) for every point $x \in M$ we have 
    \begin{equation*}
        \limsup_{r \to 0}\sup_{t \in [0, 1]}\haus^n(\Gamma_t\cap B_r(x)) = 0.
    \end{equation*}
\end{itemize}
\end{definition}
\begin{definition}[Sweepout of an open set]\label{def:sweepout}
Let $U \subset M$ be a relatively open set of $M$, where $M$ is a complete manifold with boundary. $\{\Gamma_t\}_{t \in [0,1]}$ is a sweepout of $U$ if it satisfies $(s1)-(s4)$ and there exists a family $\br{\Omega_t}$, $t \in [0, 1]$, of relative open sets of finite Hausdorff meausre, such that:
\begin{enumerate}
    \item[(sw1)] $(\Gamma_t \backslash\partial\Omega_t) \subset P_t$ for any $t$;
    \item[(sw2)] $\Omega_0\cap U = \varnothing$ and $U \subset \Omega_1$;
    \item[(sw3)] $\haus^{n + 1}(\Omega_t\backslash \Omega_2) + \haus^{n + 1}(\Omega_s\backslash\Omega_t) \to 0$ as $t \to s$
\end{enumerate}
For a sweepout $\br{\Gamma_t}$ we will say that $\br{\Omega_t}$ is the \emph{corresponding family} of (relative) open sets if it satisfies $(sw1)-(sw3)$.
\end{definition}

For a sweepout $\{\Gamma_t\}$, we define the \emph{height} or \emph{width} of the sweepout as the maximal area achieved along the family:

\begin{definition}[Height of a sweepout]\label{def:height}
The \emph{height} (or \emph{width}) of a sweepout $\{\Gamma_t\}_{t \in [0,1]}$ is defined by
\begin{equation*}
    \mathcal{F}(\{\Gamma_t\}) := \sup_{t \in [0,1]} \mathcal{H}^n(\Gamma_t\backslash\partial M).
\end{equation*}
\end{definition}


For a set $\mathcal{S}$ of sweepouts, we define:
\begin{definition}[Width of a collection]\label{def:width-collection}
Let $\mathcal{S}$ be a collection of sweepouts of $U$. The \emph{width} of $\mathcal{S}$ is
\begin{equation*}
    W(\mathcal{S}) := \inf_{\{\Gamma_t\} \in \mathcal{S}} \mathcal{F}(\{\Gamma_t\}).
\end{equation*}
This is the infimal height over all sweepouts in the collection $\mathcal{S}$.
\end{definition}

\begin{definition}[Good sweepout]\label{def:good-sweepout}
Let $U \subset M$ be a relative open subset. A sweepout $\{\Gamma_t\}$ of $U$ is
called a \emph{good sweepout} if it in addition satisfies
\begin{enumerate}
    \item[($\text{sw}_g$)] $\haus^n(\Gamma_0) \leq 5\haus^n(\partial U)$ and $\haus^n(\Gamma_1) \leq 5\haus^n(\partial U)$.
\end{enumerate}
we will denote the collection of good sweepout of $U$ as $\mathcal{S}_g(U)$ or simply $\mathcal{S}_g$ in short.
\end{definition}
\begin{definition}[Nested sweepout]\label{def:nested-sweepout}
A \emph{nested sweepout} $\br{\Gamma_t}$ is a sweepout of $U$ which in addition satisfies:
\begin{enumerate}
    \item[($\text{sw}_n$)] there exists a Morse function $f: \widetilde{M} \to [-1, \infty)$, such that $\Gamma_t = f^{-1}(t)\cap M$, $t \in [0, 1]$; the corresponding family of open sets is given by $\Omega_t = f^{-1}((-\infty, t))\cap M$.
\end{enumerate}
we will denote the collection of nested sweepout of $U$ as $\mathcal{S}_n(U)$ or simply $\mathcal{S}_n$ in short.
\end{definition}

\begin{definition}[Relative sweepout]\label{def:relative-sweepout}
Suppose $\partial U$ is a smooth manifold and $\br{\Gamma_t}$ is a nested sweepout of $U$ with the corresponding open sets $\br{\Omega_t}$. Set $\Sigma_t = (cl(U)\cap \Gamma_t)$. We will say that $\br{\Sigma_t}$ is a \emph{relative sweepout} of $U$. we will denote the collection of relative sweepout of $U$ as $\mathcal{S}_\partial(U)$ or simply $\mathcal{S}_\partial$ in short.
\end{definition}

With respect to different family of sweepouts, we can define their corresponding width $W_n(U)$, $W_g(U)$ and $W_\partial(U)$.
These width functionals satisfy the following relationships~\cite{CL20}:
\begin{proposition}[Width inequalities]\label{prop:width-inequalities}
Let $U \subset M$ be a bounded relatively open set. Then:
\begin{enumerate}[label=(\roman*)]
    \item $W_\partial(U) \leq W_n(U) \leq W_\partial(U) + \mathcal{H}^n(\bdy U)$;
    \item $W_n(U) = W(U)$;
    \item If $U$ is a good set, then $W_g(U) = W(U)$.
\end{enumerate}
\end{proposition}

\begin{definition}[Good set]\label{def:good-set}
A bounded relative open subset $U \subset M$ (or equivalently, a bounded Caccioppoli set $E$ with
$U = \mathrm{int}(E)$) is a \emph{good set} if
\[
    \mathcal{H}^n(\bdy U) \leq \frac{1}{10}\,W_\partial(U).
\]
That is, the $n$-dimensional Hausdorff measure of the boundary $\bdy U$ is at most $1/10$ of the relative
width.
\end{definition}
\begin{remark}[Why good sweepouts prevent escape to infinity]\label{rem:good-sweepout-no-escape}
The combination of the good set condition $\mathcal{H}^n(\bdy U) \leq \frac{1}{10} W_\partial(U)$ and the good sweepout condition $\mathcal{F}(\{\Gamma_t\}) \leq 5 W_\partial(U)$ is designed to address the fundamental difficulty of \emph{escape to infinity} in min-max theory on complete non-compact manifolds.

\textbf{The escape to infinity problem:} In a compact manifold, the min-max surface is automatically localized by compactness. However, in a complete non-compact manifold (such as $\mathbb{R}^{n+1}$ or a manifold with cylindrical ends), a sequence of surfaces in a sweepout could concentrate all their mass arbitrarily far from the region $U$ where we want to find the minimal surface, thus ``escaping to infinity.''

\textbf{How good sweepouts prevent escape:} The key observation is the following \emph{bottleneck principle}:
\begin{quote}
If a sweepout wants to transfer mass from inside $U$ to outside $U$, it must pass through the boundary $\bdy U$. But $\bdy U$ has area $\leq W_\partial(U)/10$, which is much smaller than the sweepout's height $\sim W_\partial(U)$. Therefore, the sweepout cannot move all its mass outside $U$---at least one slice must carry a definite amount of area inside $\overline{U}$.
\end{quote}
This principle is made rigorous in Chapter~\ref{ch:no-escape-to-infinity} (Proposition~\ref{prop:non-escape}), which proves that every good sweepout must have at least one slice with area $\geq \varepsilon(U) > 0$ inside $\overline{U}$. The proof uses a contradiction argument: if all slices had negligible area in $U$, we could use the nested sweepout structure (Chapter~\ref{ch:nested-sweepout}) to ``push'' them completely outside $U$ while increasing the perimeter by at most $\mathcal{H}^n(\bdy U) \ll W_\partial(U)$, thus constructing a competing sweepout with strictly smaller width---a contradiction.
\end{remark}

\begin{theorem}[Existence of good sweepouts]\label{thm:existence-good-sweepouts}
Let $U \subset M$ be a good set. Then there exists a good sweepout $\{\Gamma_t\}$ of $U$, i.e., a sweepout satisfying $\mathcal{F}(\{\Gamma_t\}) \leq 5 W_\partial(U)$.
\end{theorem}

\begin{theorem}[Improvement of sweepouts]\label{thm:improvement-sweepouts}
Let $U \subset M$ be a good set, and let $\{\Gamma_t\}$ be a sweepout of $U$ with
$\mathcal{F}(\{\Gamma_t\}) < W_\partial(U)$. Then $\{\Gamma_t\}$ can be homotoped to a sweepout
$\{\Gamma_t'\}$ with strictly smaller height: $\mathcal{F}(\{\Gamma_t'\}) < \mathcal{F}(\{\Gamma_t\})$.
\end{theorem}